
\documentclass[addpoints]{exam}
\usepackage{amsmath}
\usepackage{amssymb}
\usepackage{color}
\usepackage[version=4]{mhchem}
\usepackage{siunitx}

% \printanswers

\newcommand{\event}{Final Exam}
\newcommand{\tournament}{}
\def\type{0}

\setlength\fillinlinelength{\textwidth}
\CorrectChoiceEmphasis{\color{red}\bfseries}
\SolutionEmphasis{\color{red}\bfseries}

\setlength\answerclearance{2pt}
\pagestyle{headandfoot}
\runningheadrule
\runningheader{SCH3U \event}{\tournament}{\ifnum\type=1 Class Set Test \else Full Name:\hspace{1cm} \fi}
\runningfootrule
\runningfooter{}{Page \thepage\ of \numpages}{}

\begin{document}
\begin{center}
    {\huge Grade 11 Chemistry \\ \event
    \ifprintanswers \space Key
    \else \space \ifnum\type<2 \fi
    \fi \\ \vspace{3mm}\tournament\vspace{1mm}}
\end{center}

\vspace{.20\textheight}

\begin{center}
    First Name: \ifprintanswers
    \underline{\hspace{3.5cm}}\textbf{KEY}\underline{\hspace{5.5cm}}\\\vspace{5mm}
    \else \underline{\hspace{10cm}}\\ \vspace{5mm} \fi
    Last Name: \underline{\hspace{10cm}}\\ \vspace{5mm}
\end{center}

\noindent\textbf{\underline{Directions:}}
\begin{itemize}
    \item Show all work with units. A periodic table is provided.
    \item Use $N_A = 6.022\times10^{23}$\,mol$^{-1}$;\;
    $R = \SI{8.314}{kPa\cdot L/(mol\cdot K)}$;\;
    Molar volume at STP $= \SI{22.4}{L/mol}$.
    \item This exam covers all units of SCH3U. Designed for 2 hours.
\end{itemize}
\begin{center}
\textit{For grading use only}\\
\multirowgradetable{2}[pages]
\end{center}

\newpage
\begin{questions}

\section*{Part A --- Multiple Choice (20 marks)}
\vspace{3mm}
\noindent\textit{Unit 1: Matter, Bonding \& Trends}
\vspace{3mm}

\question[1] Which set correctly ranks \textbf{atomic radius} from largest to smallest?
\begin{choices}
    \choice $\text{F} > \text{Cl} > \text{Br}$
    \correctchoice $\text{K} > \text{Na} > \text{Li}$
    \choice $\text{O} > \text{S} > \text{Se}$
    \choice $\text{Na} > \text{Mg} > \text{Al}$ and $\text{Al} > \text{Na}$
\end{choices}

\question[1] The molecular geometry of \ce{NH3} is:
\begin{choices}
    \choice Tetrahedral
    \correctchoice Trigonal pyramidal
    \choice Trigonal planar
    \choice Bent
\end{choices}

\question[1] Which species has the \textbf{same} electron configuration as \ce{Ne}?
\begin{choices}
    \choice \ce{Na+} only
    \choice \ce{F-} only
    \correctchoice Both \ce{Na+} and \ce{F-}
    \choice \ce{Mg^{2+}} only
\end{choices}

\question[1] Which molecule has a \textbf{net dipole moment}?
\begin{choices}
    \choice \ce{CCl4}
    \choice \ce{CO2}
    \correctchoice \ce{H2O}
    \choice \ce{BF3}
\end{choices}

\vspace{3mm}
\noindent\textit{Unit 2: Chemical Reactions}
\vspace{3mm}

\question[1] In the single displacement reaction \ce{Mg + 2HCl -> MgCl2 + H2},
the oxidizing agent is:
\begin{choices}
    \choice Mg
    \correctchoice \ce{HCl} (specifically \ce{H+})
    \choice \ce{MgCl2}
    \choice \ce{H2}
\end{choices}

\question[1] Which product forms when solutions of \ce{Pb(NO3)2} and \ce{KI} are mixed?
\begin{choices}
    \choice \ce{KNO3} precipitate
    \correctchoice \ce{PbI2} precipitate
    \choice \ce{Pb(OH)2} precipitate
    \choice No precipitate forms
\end{choices}

\question[1] The balanced equation for complete combustion of \ce{C4H10} is:
\begin{choices}
    \choice \ce{C4H10 + 5O2 -> 4CO2 + 5H2O}
    \correctchoice \ce{2C4H10 + 13O2 -> 8CO2 + 10H2O}
    \choice \ce{C4H10 + 6O2 -> 4CO2 + 5H2O}
    \choice \ce{C4H10 + 13O2 -> 4CO2 + 5H2O}
\end{choices}

\question[1] The net ionic equation for any strong acid--strong base neutralization is:
\begin{choices}
    \choice \ce{H2O -> H+ + OH-}
    \correctchoice \ce{H+(aq) + OH-(aq) -> H2O(l)}
    \choice \ce{HA + BOH -> BA + H2O}
    \choice \ce{H+ + B+ -> HB+}
\end{choices}

\vspace{3mm}
\noindent\textit{Unit 3: Stoichiometry}
\vspace{3mm}

\question[1] How many grams of \ce{H2O} are produced when $1.00$ mol of \ce{C3H8}
undergoes complete combustion? (\ce{C3H8 + 5O2 -> 3CO2 + 4H2O})
\begin{choices}
    \choice \SI{18.0}{g}
    \correctchoice \SI{72.0}{g}
    \choice \SI{36.0}{g}
    \choice \SI{54.0}{g}
\end{choices}

\question[1] A compound is 75.0\% C and 25.0\% H by mass. Its empirical formula is:
\begin{choices}
    \choice \ce{CH}
    \correctchoice \ce{CH4}
    \choice \ce{C2H6}
    \choice \ce{CH3}
\end{choices}

\question[1] The limiting reagent in a reaction is:
\begin{choices}
    \choice The reagent with the greatest mass
    \correctchoice The reagent that produces the least product
    \choice The reagent with the smallest molar mass
    \choice The reagent present in the smallest number of moles
\end{choices}

\question[1] The percent yield of a reaction is 85\%. If the actual yield is $\SI{34}{g}$,
the theoretical yield is:
\begin{choices}
    \choice \SI{29}{g}
    \correctchoice \SI{40}{g}
    \choice \SI{39}{g}
    \choice \SI{34}{g}
\end{choices}

\vspace{3mm}
\noindent\textit{Unit 4: Solutions}
\vspace{3mm}

\question[1] What volume of $\SI{0.500}{mol/L}$ \ce{KOH} contains $0.125$ mol of \ce{KOH}?
\begin{choices}
    \choice \SI{62.5}{mL}
    \correctchoice \SI{250}{mL}
    \choice \SI{400}{mL}
    \choice \SI{500}{mL}
\end{choices}

\vspace{3mm}
\noindent\textit{Unit 5: Gases}
\vspace{3mm}

\question[1] A gas at $\SI{300}{K}$ and $\SI{200}{kPa}$ occupies $\SI{5.0}{L}$.
At $\SI{600}{K}$ and $\SI{100}{kPa}$, its volume is:
\begin{choices}
    \choice \SI{5.0}{L}
    \choice \SI{10.0}{L}
    \correctchoice \SI{20.0}{L}
    \choice \SI{2.5}{L}
\end{choices}

\question[1] At STP, the density of \ce{O2} gas ($M = \SI{32.0}{g/mol}$) is:
\begin{choices}
    \choice \SI{1.00}{g/L}
    \correctchoice \SI{1.43}{g/L}
    \choice \SI{2.00}{g/L}
    \choice \SI{22.4}{g/L}
\end{choices}

\question[1] In a mixture of gases, each gas exerts a pressure as if it were the only
gas present. This is:
\begin{choices}
    \choice Boyle's Law
    \choice Charles's Law
    \correctchoice Dalton's Law of Partial Pressures
    \choice The Ideal Gas Law
\end{choices}

\question[1] How many moles of gas are in \SI{5.60}{L} at STP?
\begin{choices}
    \choice $0.100$ mol
    \correctchoice $0.250$ mol
    \choice $0.500$ mol
    \choice $1.00$ mol
\end{choices}

\question[1] A gas occupies $\SI{3.0}{L}$ at $\SI{27}{\celsius}$. At what temperature
(in \si{\celsius}) will it occupy $\SI{4.5}{L}$ at constant pressure?
\begin{choices}
    \choice $\SI{40.5}{\celsius}$
    \correctchoice $\SI{177}{\celsius}$
    \choice $\SI{450}{\celsius}$
    \choice $\SI{150}{\celsius}$
\end{choices}

\question[1] The ideal gas law assumes:
\begin{choices}
    \choice Gas particles have significant volume
    \correctchoice No intermolecular forces between gas particles
    \choice Gas particles move slowly at high temperatures
    \choice Collisions between particles are inelastic
\end{choices}

\question[1] A container holds \ce{N2} at \SI{50}{kPa}, \ce{O2} at \SI{30}{kPa}, and
\ce{CO2} at \SI{20}{kPa}. The mole fraction of \ce{N2} is:
\begin{choices}
    \choice $0.3$
    \correctchoice $0.5$
    \choice $0.6$
    \choice $1.0$
\end{choices}


\section*{Part B --- Short Answer (40 marks)}

\question \textbf{(Bonding \& Trends)} Answer the following.
\begin{parts}
    \part[2] Draw the Lewis structure of \ce{SO2}. Include all lone pairs, indicate
    bond type, and predict the molecular geometry.
    \part[3] Arrange \ce{HF}, \ce{HCl}, \ce{HBr}, \ce{HI} in order of increasing
    boiling point. Explain the trend, accounting for the anomalously high boiling
    point of \ce{HF}.
\end{parts}
\begin{solution}[9cm]
\textbf{(a)} \ce{SO2}: S is central with one double bond to O and one single bond with a
formal charge, plus 1 lone pair on S. Total electron geometry: bent (trigonal planar electron
geometry). \textbf{Molecular geometry: bent}. Bond angle $\approx 119^\circ$.
(\ce{SO2} has resonance structures; both S--O bonds are equivalent.)

\textbf{(b)} Order of increasing boiling point: \ce{HCl} $<$ \ce{HBr} $<$ \ce{HI} $<$ \ce{HF}.

For \ce{HCl} through \ce{HI}: boiling point increases with molar mass due to stronger London
dispersion forces. \ce{HF} has an anomalously \textit{high} boiling point despite its low molar
mass because \ce{HF} molecules form strong \textbf{hydrogen bonds} (F is the most
electronegative element, making the H--F bond highly polar).
\end{solution}

\question \textbf{(Reactions)} For each, write a balanced molecular equation, identify
the reaction type, and (where applicable) write the net ionic equation.
\begin{parts}
    \part[3] Aqueous solutions of iron(III) chloride and sodium hydroxide are mixed.
    \part[3] Solid calcium carbonate is heated strongly.
\end{parts}
\begin{solution}[9cm]
\textbf{(a)} \ce{Fe(OH)3} is insoluble (precipitate forms). Type: \textbf{double displacement}.
\[
\ce{FeCl3(aq) + 3NaOH(aq) -> Fe(OH)3(s) + 3NaCl(aq)}
\]
Net ionic: \ce{Fe^{3+}(aq) + 3OH^{-}(aq) -> Fe(OH)3(s)}

\textbf{(b)} Type: \textbf{decomposition}.
\[
\ce{CaCO3(s) ->[\Delta] CaO(s) + CO2(g)}
\]
(No net ionic equation — no aqueous ions involved.)
\end{solution}

\question \textbf{(Stoichiometry)} Aluminum reacts with oxygen:
\[\ce{4Al(s) + 3O2(g) -> 2Al2O3(s)}\]
$M(\ce{Al}) = \SI{27.0}{g/mol}$;\quad $M(\ce{Al2O3}) = \SI{102}{g/mol}$
\begin{parts}
    \part[2] How many grams of \ce{Al2O3} are produced from $\SI{13.5}{g}$ of Al
    with excess \ce{O2}?
    \part[3] If $\SI{13.5}{g}$ of Al and $\SI{12.0}{g}$ of \ce{O2} are used, identify
    the limiting reagent and find the theoretical yield of \ce{Al2O3}.
    \part[2] A student obtains $\SI{22.0}{g}$ of \ce{Al2O3}. Calculate the percent yield.
\end{parts}
\begin{solution}[9cm]
\textbf{(a)} $n(\ce{Al}) = \dfrac{13.5}{27.0} = 0.500$ mol;\quad
$n(\ce{Al2O3}) = 0.500 \times \dfrac{2}{4} = 0.250$ mol;\quad
$m = 0.250 \times 102 = \mathbf{\SI{25.5}{g}}$

\textbf{(b)} $n(\ce{O2}) = \dfrac{12.0}{32.0} = 0.375$ mol

\ce{Al2O3} from Al: $0.250$ mol;\quad \ce{Al2O3} from \ce{O2}: $0.375 \times \dfrac{2}{3} = 0.250$ mol

Both give $0.250$ mol — they are in \textbf{exact stoichiometric ratio} (neither is limiting).
Theoretical yield $= \mathbf{\SI{25.5}{g}}$

\textbf{(c)} $\%\text{ yield} = \dfrac{22.0}{25.5}\times100 = \mathbf{86.3\%}$
\end{solution}

\question \textbf{(Solutions)} A student dissolves $\SI{5.85}{g}$ of \ce{NaCl}
($M = \SI{58.5}{g/mol}$) to make $\SI{500.0}{mL}$ of solution.
\begin{parts}
    \part[2] Calculate the concentration of the \ce{NaCl} solution.
    \part[2] Calculate $[\ce{Na+}]$ and $[\ce{Cl-}]$ in the solution.
    \part[3] This solution is then diluted: $\SI{75.0}{mL}$ is taken and diluted to
    $\SI{300.0}{mL}$. Calculate the concentration of \ce{NaCl} in the diluted solution.
\end{parts}
\begin{solution}[8cm]
\textbf{(a)} $n = \dfrac{5.85}{58.5} = 0.100$ mol;\quad
$C = \dfrac{0.100}{0.5000} = \mathbf{\SI{0.200}{mol/L}}$

\textbf{(b)} \ce{NaCl -> Na+ + Cl-};\quad $[\ce{Na+}] = [\ce{Cl-}] = \mathbf{\SI{0.200}{mol/L}}$

\textbf{(c)} $C_1V_1 = C_2V_2$:\quad
$C_2 = \dfrac{0.200 \times 75.0}{300.0} = \mathbf{\SI{0.0500}{mol/L}}$
\end{solution}

\question \textbf{(Gases)} A \SI{2.50}{L} cylinder contains a gas at $\SI{25}{\celsius}$
and $\SI{500}{kPa}$.
\begin{parts}
    \part[3] Calculate the number of moles of gas in the cylinder.
    \part[3] The cylinder is heated to $\SI{100}{\celsius}$ at constant volume. What is the
    new pressure?
    \part[2] The gas is identified as \ce{CO2} ($M = \SI{44.0}{g/mol}$). What mass of
    gas is in the cylinder?
\end{parts}
\begin{solution}[9cm]
\textbf{(a)}
\[
n = \frac{PV}{RT} = \frac{500 \times 2.50}{8.314 \times 298} = \frac{1250}{2477.6}
= \mathbf{0.5046 \text{ mol}}
\]

\textbf{(b)} $T_1 = 298\,\text{K}$,\; $T_2 = 373\,\text{K}$.
\[
P_2 = P_1 \cdot \frac{T_2}{T_1} = 500 \times \frac{373}{298} = \mathbf{\SI{625.8}{kPa}}
\]

\textbf{(c)} $m = n \times M = 0.5046 \times 44.0 = \mathbf{\SI{22.2}{g}}$
\end{solution}

\end{questions}
\end{document}
